\documentclass[10pt,a4paper]{amsart}
\usepackage{amsmath,amssymb,amsthm,graphicx}
% For commutative diagrams we use Xy-pic
\usepackage[all]{xy}

% We use 2cell for 2-commutative diagrams.
\xyoption{2cell}
% Theorem environments.
%
\theoremstyle{plain}
\newtheorem{theorem}[subsection]{Theorem}
\newtheorem{proposition}[subsection]{Proposition}
\newtheorem{lemma}[subsection]{Lemma}

\theoremstyle{definition}
\newtheorem{definition}[subsection]{Definition}
\newtheorem{example}[subsection]{Example}
\newtheorem{exercise}[subsection]{Exercise}
\newtheorem{situation}[subsection]{Situation}

\theoremstyle{remark}
\newtheorem{remark}[subsection]{Remark}
\newtheorem{remarks}[subsection]{Remarks}

\numberwithin{equation}{subsection}

% Macros
%
\def\lim{\mathop{\mathrm{lim}}\nolimits}
\def\colim{\mathop{\mathrm{colim}}\nolimits}
\def\Spec{\mathop{\mathrm{Spec}}}
\def\Hom{\mathop{\mathrm{Hom}}\nolimits}
\def\Ext{\mathop{\mathrm{Ext}}\nolimits}
\def\SheafHom{\mathop{\mathcal{H}\!\mathit{om}}\nolimits}
\def\SheafExt{\mathop{\mathcal{E}\!\mathit{xt}}\nolimits}
\def\Sch{\mathit{Sch}}
\def\Mor{\mathop{\mathrm{Mor}}\nolimits}
\def\Ob{\mathop{\mathrm{Ob}}\nolimits}
\def\Sh{\mathop{\mathit{Sh}}\nolimits}
\def\NL{\mathop{N\!L}\nolimits}
\def\CH{\mathop{\mathrm{CH}}\nolimits}
\def\proetale{{pro\text{-}\acute{e}tale}}
\def\etale{{\acute{e}tale}}
\def\QCoh{\mathit{QCoh}}
\def\Ker{\mathop{\mathrm{Ker}}}
\def\Im{\mathop{\mathrm{Im}}}
\def\Coker{\mathop{\mathrm{Coker}}}
\def\Coim{\mathop{\mathrm{Coim}}}

% Boxtimes
%
\DeclareMathSymbol{\boxtimes}{\mathbin}{AMSa}{"02}

%
% Macros for moduli stacks/spaces
%
\def\QCohstack{\mathcal{QC}\!\mathit{oh}}
\def\Cohstack{\mathcal{C}\!\mathit{oh}}
\def\Spacesstack{\mathcal{S}\!\mathit{paces}}
\def\Quotfunctor{\mathrm{Quot}}
\def\Hilbfunctor{\mathrm{Hilb}}
\def\Curvesstack{\mathcal{C}\!\mathit{urves}}
\def\Polarizedstack{\mathcal{P}\!\mathit{olarized}}
\def\Complexesstack{\mathcal{C}\!\mathit{omplexes}}
% \Pic is the operator that assigns to X its picard group, usage \Pic(X)
% \Picardstack_{X/B} denotes the Picard stack of X over B
% \Picardfunctor_{X/B} denotes the Picard functor of X over B
\def\Pic{\mathop{\mathrm{Pic}}\nolimits}
\def\Picardstack{\mathcal{P}\!\mathit{ic}}
\def\Picardfunctor{\mathrm{Pic}}
\def\Deformationcategory{\mathcal{D}\!\mathit{ef}}
\newenvironment{reference}{}{}
\newenvironment{slogan}{}{}
\newenvironment{history}{}{}
\usepackage{cleveref}
\crefname{originallemma}{original lemma}{original lemmas}
\Crefname{originallemma}{Original lemma}{Original lemmas}
\makeatletter
\newlabel{lemma-divisor-special-fiber}{{0C5Z}{source}{}{sourcepointer-1}{}}
\newlabel{lemma-divisor-special-fiber@cref}{{[originallemma][0][]0C5Z}{[1][0][]source}{}{}{}}
\makeatother
\crefname{originallemma}{original lemma}{original lemmas}
\Crefname{originallemma}{Original lemma}{Original lemmas}
\makeatletter
\newlabel{lemma-regular-model-pic}{{0C63}{source}{}{sourcepointer-2}{}}
\newlabel{lemma-regular-model-pic@cref}{{[originallemma][0][]0C63}{[1][0][]source}{}{}{}}
\makeatother
\crefname{originallemma}{original lemma}{original lemmas}
\Crefname{originallemma}{Original lemma}{Original lemmas}
\makeatletter
\newlabel{lemma-canonical-map-of-pic}{{0CAC}{source}{}{sourcepointer-3}{}}
\newlabel{lemma-canonical-map-of-pic@cref}{{[originallemma][0][]0CAC}{[1][0][]source}{}{}{}}
\makeatother
\crefname{originallemma}{original lemma}{original lemmas}
\Crefname{originallemma}{Original lemma}{Original lemmas}
\makeatletter
\newlabel{lemma-characterize-trivial}{{0CAB}{source}{}{sourcepointer-4}{}}
\newlabel{lemma-characterize-trivial@cref}{{[originallemma][0][]0CAB}{[1][0][]source}{}{}{}}
\makeatother
\title{Torsion of order coprime to the gcd of the fibre multiplicities gives an exact sequence of Picard groups}
\author{The Stacks Project Authors}
\date{stacks-project:models.tex:local-snapshot-20261012; GNU FDL 1.2 or later}
\begin{document}
\maketitle
\noindent\textit{Editorial note.} The counted unit is the Stacks Project Lemma lemma-sequence-torsion at source lines 5962-5970 together with its proof at source lines 5972-5992, a straight adjacent pairing in the chapter file. The chapter is the upstream file models.tex, which is the file name recorded in the source record; the pool directory carries the longer name stacks-models and is not used as the chapter name anywhere in this card.\par\medskip
\noindent\textit{Editorial note.} The retained statement opens with the words In Situation followed by a reference to situation-regular-model, so the situation environment that defines the setting is retained as a context span at source lines 4628-4637. It is context, not a counted unit, and it does not overlap the counted interval.\par\medskip
\noindent\textit{Editorial note.} This chapter's candidate index records the statement body and the proof body without their closing markers, so the retained spans used here are the full environments, exactly as in the earlier cards of this chapter.\par\medskip
\noindent\textit{Editorial note.} The retained proof names three results of the same chapter, for the exact sequence of the regular model, for the canonical map to the combinatorial Picard group and for the characterisation of the triviality of a power of an invertible sheaf; all three are declared as non-load-bearing navigation pointers with their stable Stacks tags looked up in the shared tag index. The retained situation additionally names one lemma of this chapter, also declared as a pointer.\par\medskip
\noindent\textit{Editorial note.} No citation occurs in the retained spans, so no bibliography entry is supplied, and the retained preamble is the chapter preamble's Xy-pic setup at lines 14-18, the theorem declarations and mathematical macros at lines 154-217, and the three no-op wrapper environments.\par\medskip
\noindent\textit{Allowed standard prerequisites (editorial).}\begin{itemize}
\item Exact sequences of Picard groups for a regular model and torsion subgroups. For a regular model the Picard group of the model surjects onto that of the curve with kernel determined by the multiplicities of the special fibre, and there is a canonical map from the Picard group of the curve to the combinatorial Picard group of the numerical type whose kernel consists of the restrictions of invertible sheaves of degree zero on every component; an invertible sheaf whose restriction to each component has degree zero becomes torsion of order dividing the gcd of the multiplicities after raising to a suitable power. The author's own lemmas of this section are the allowed inputs.. The counted sequence is the tool used in this section to compare the h-torsion of the three Picard groups for an order h coprime to the gcd of the multiplicities, which is the shape needed for the semistable-reduction results.
\end{itemize}\medskip
\ifdefined\hypertarget\hypertarget{sourcepointer-1}{}\fi
\noindent\textit{Source reference.} Original lemma 0C5Z: focused-seeds/AGP/sources/stacks-snapshot-20261011/tags.txt, source lines 15279--15279. Navigation pointer inside the retained situation, which names it when it writes the special fibre as a sum of its components with multiplicities; the gcd of those multiplicities is part of the statement. The number is the stable Stacks tag of the target, looked up in the shared tag index.\par\medskip
\ifdefined\hypertarget\hypertarget{sourcepointer-2}{}\fi
\noindent\textit{Source reference.} Original lemma 0C63: focused-seeds/AGP/sources/stacks-snapshot-20261011/tags.txt, source lines 15283--15283. Pointer to the exact sequence of the regular model, which the retained proof takes h-torsion of; that exact sequence is counted by a separate card of this lane and is not reproduced here. The number is the stable Stacks tag of the target, looked up in the shared tag index.\par\medskip
\ifdefined\hypertarget\hypertarget{sourcepointer-3}{}\fi
\noindent\textit{Source reference.} Original lemma 0CAC: focused-seeds/AGP/sources/stacks-snapshot-20261011/tags.txt, source lines 15432--15432. Pointer to the canonical map from the Picard group of the curve to that of the numerical type, which the retained proof uses; that map is counted by a separate card of this lane and is not reproduced here. The number is the stable Stacks tag of the target, looked up in the shared tag index.\par\medskip
\ifdefined\hypertarget\hypertarget{sourcepointer-4}{}\fi
\noindent\textit{Source reference.} Original lemma 0CAB: focused-seeds/AGP/sources/stacks-snapshot-20261011/tags.txt, source lines 15431--15431. Pointer to the characterisation of the invertible sheaves whose h-th tensor power is trivial, used by the retained proof for the order of that power; the target is an allowed input of the same section and is not reproduced here. The number is the stable Stacks tag of the target, looked up in the shared tag index.\par\medskip
\begin{situation}
\label{situation-regular-model}
Let $R$ be a discrete valuation ring with fraction field $K$,
residue field $k$, and uniformizer $\pi$.
Let $C$ be a smooth projective curve over $K$ with $H^0(C, \mathcal{O}_C) = K$.
Let $X$ be a regular proper model of $C$.
Let $C_1, \ldots, C_n$ be the irreducible components of the special
fibre $X_k$. Write $X_k = \sum m_i C_i$ as in
Lemma \ref{lemma-divisor-special-fiber}.
\end{situation}

\begin{lemma}
\label{lemma-sequence-torsion}
In Situation \ref{situation-regular-model} let $d = \gcd(m_1, \ldots, m_n)$
and let $T$ be the numerical type associated to $X$.
Let $h \geq 1$ be an integer prime to $d$. There exists an exact sequence
$$
0 \to \Pic(X)[h] \to \Pic(C)[h] \to \Pic(T)[h]
$$
\end{lemma}

\begin{proof}
Taking $h$-torsion in the exact sequence of
Lemma \ref{lemma-regular-model-pic}
we obtain the exactness of
$0 \to \Pic(X)[h] \to \Pic(C)[h]$
because $h$ is prime to $d$.
Using the map of Lemma \ref{lemma-canonical-map-of-pic}
we get a map $\Pic(C)[h] \to \Pic(T)[h]$
which annihilates elements of $\Pic(X)[h]$.
Conversely, if $\xi \in \Pic(C)[h]$
maps to zero in $\Pic(T)[h]$, then we can find
an invertible $\mathcal{O}_X$-module $\mathcal{L}$
with $\deg(\mathcal{L}|_{C_i}) = 0$ for all $i$
whose restriction to $C$ is $\xi$.
Then $\mathcal{L}^{\otimes h}$ is $d$-torsion by
Lemma \ref{lemma-characterize-trivial}.
Let $d'$ be an integer such that $dd' \equiv 1 \bmod h$.
Such an integer exists because $h$ and $d$ are coprime.
Then $\mathcal{L}^{\otimes dd'}$ is an $h$-torsion
invertible sheaf on $X$ whose restriction to $C$ is $\xi$.
\end{proof}

\end{document}
